\chapter[常存质疑]{常存质疑}

\chaptersubtitle{《奇迹、魔法，都是存在的》}

\chapternote{%
    ``I think theatre should \textbf{always be somewhat suspect}.''
    \par
    {\raggedleft\noindent ---Václav Havel, \emph{Disturbing the Peace}\par}
    \par
    “戏剧应常存质疑。”
    \par
    {\raggedleft\noindent ——瓦茨拉夫·哈维尔《扰乱和平》\par}
    \par
    译者按：第四集从麻美身后的空洞中浮现的，正是这种不肯把答案当作答案的怀疑。
}

《魔法少女小圆》第四集开启了该剧的第二个篇章。
然而，此篇章的核心焦点——成为魔法少女的沙耶香与新角色佐仓杏子——直至集末才真正展开。
本集主体实则起着承上启下的作用，既充当了第一篇章的尾声，又为第二篇章埋下引线。

此外，本集可以说才是动画“真正的”第一集。
相较于前三集对传统魔法少女风格的表面效仿，本集首次鲜明地呈现了作品的主导美学风格。
因此，它亦可视为全剧的正式引介：
作品由此首度揭示了贯穿动画的核心主题之一——
抑郁。

在论及《魔圆》时，人们经常引用编剧虚渊玄的一段言论（尽管原文出自其早期作品 \emph{Fate/Zero} 第一卷后记\footnote{%
    本书仅有日文版。相关后记的非官方译文可见 Baka-Tsuki。
}）。
他写道：
“事实上，我并非向来如此。我也曾创作过结局未必完满，但主角直至终章仍秉持‘纵使前路艰险，亦须坚守初心’的作品。然而不知何时起，我再难创作此类故事”\cite{urobuchi2006fate-zero}。
他继而又论及熵增与失败的必然性，言辞间流露出某种抑郁迹象，至少也显得颇为消极。

诚然，依据作品揣测作者的精神状态或个人观点风险重重。
“隐含作者\footnote{%
    译注：由美国文学理论家韦恩·布斯在《小说修辞学》中提出，指文本所呈现、由读者建构的作者形象，区别于真实存在的血肉之躯。
    故即便真实作者否认，文学批评仍可将文本蕴含的价值观、要素与手法归于其隐含作者。
}”必然有别于实际作者\cite{booth1961rhetoric-fiction}。
后者作为真实个体（迥异于角色或隐含作者这类虚构存在），其主观内心世界全然不可知，无法从外部感知或推断，尽管我们仍需假定其存在。
换言之，虚渊玄本人或许生活顺遂，塑造抑郁角色与绝望情境，纯粹出于对此类题材的偏好与驾驭能力。
真相如何，我们无从知晓。

尽管如此，将《魔法少女小圆》解读为一部关于抑郁的故事依然顺理成章。
剧中的魔法少女均呈现不同的抑郁征象：
麻美深陷孤独，焰情感淡漠，沙耶香则在贯穿始终的失落与层层累积的绝望中，最终走向了自杀。

作为动画首个“真正的”开篇，本集弥漫着抑郁与失落。
前半部分着重刻画小圆如何应对失去麻美的创痛。
麻美去世后的翌晨，早餐蛋黄的色泽令她忆起麻美的金发，这滋味更是催人泪下。
此情此景，令人联想到另一部备受喜爱的魔法少女真人剧中的经典台词\footnote{%
    译注：出自《吸血鬼猎人巴菲》第五季第 16 集（约 26 分 53 秒至 27 分 30 秒）。
    主人公巴菲面对自身能力有限、无法从死亡中挽救母亲这一现实时，所说的台词。
}：

“但我不明白！
我不明白为什么事情会发生，也不明白要怎么挺过去，但她现在只是……
只剩一具躯壳了。
我不明白她为什么不能回到身体里，然后活过来，不要死去。
这很愚蠢，极端而又愚蠢。
……我喝着果汁时在想：
她再也无法喝到果汁了，也不能吃煎蛋、打哈欠或梳头发了，再也不会了。
但是没有人跟我解释为什么。”
\cite{buffy-body-2010}

然而，本集中的失落远不止于麻美之死与小圆的哀恸。
恭介丧失演奏能力后，只能聆听音乐，其痛苦与愤怒同样撼动人心。
得知医生断言恭介的手腕残疾已成定局，沙耶香毅然决定成为魔法少女以拯救他，即便深知其中凶险。
尽管后续剧情会揭示她更复杂的动机，但此刻的她无疑是无私而英勇的。
不幸的是，恭介自残手臂时的麻木感已成预兆：
沙耶香很快也将堕入一种麻木的绝望，这种状态让她感受不到痛苦，也隔绝了一切其他感受。
此外，她选择成为麻美的继任者——
即在未及处理自身悲伤与失落前，便贸然投身于拯救他人——
也预示了两人共有的行为模式。

当然，还有那个自杀团体。
仁美在剧中戏份寥寥，她的加入显得十分突兀；
但剧情暗示，正是魔女将脆弱的个体聚拢成这样一个团体。
究竟是何使仁美变得如此脆弱？
剧中并未明言；
第二集中那位无名女子的遭遇也同样成谜。
但仍有迹可循：
此前仁美被刻画为极度忙碌：
不仅要参加优等生必需的备考补习，还要接受日本上流阶层的文化训练，如传统舞蹈或茶道。
本集之后，这些课程再无提及，或许暗示她在事件后减轻了负担。
另一种可能与她对恭介的暗恋和羞怯有关。
她在本集之后不久向沙耶香发出最后通牒；
在此之前，就连闺蜜也不知道这份感情，随后她便主动接近了恭介。

让我们把焦点更多地集中在小圆的脆弱性上：
魔女利用她未能拯救麻美的自责来发动攻击。
结合自杀团体中那位自述经营工厂失败的男人以及仁美可能的脆弱点来看，此魔女的作案模式似乎是针对目标的失败感与无能感。

至少对于小圆而言，魔女选择的攻击手段是：
通过电视屏幕播放上一集的片段，迫使她直面自己的无能。
此时小圆的形象开始模糊扭曲，用来区分角色与背景的轮廓线逐渐消融。
换言之，魔女向小圆揭示了她身为剧中角色的本质，剥夺了她的身份。
她不再是人，而是一件物品，是动画中的一个元素，因此开始消解，融入动画的背景之中。

然而，一个有趣的事实是：
纵然角色是虚构的，她们所激发的情感却是真实的。
即便知晓小圆是虚构的，我们依然会关心她，并为沙耶香前来救援感到欣慰。
至少此刻，这足以让她在剧中的“现实”里恢复完整。
在动画的领域内，小圆仍是一个独立的实体，而非弥散的概念。
尽管如此，本集作为“真正的”《魔圆》，已为全剧结局埋下了伏笔：
小圆终将化作星辰，散于天际。

至本集尾声，第二篇章已全然展开：
沙耶香成为魔法少女，其战斗服使她宛如英勇的守护骑士；
她誓言要使恭介、小圆以及世界免受魔女之害。
不幸的是，与她针锋相对的是其完美的对立面——
经验老道对应天真懵懂，炽烈红色对应沉静蓝色，自我中心的贪婪对应自我牺牲的守护。
至少表面上是如此，但与剧中的众多对立组合一样，她们终究并非那般不同。
